% =============================================================================
% TAILORED CV — Chef de Produit Téléphonie | Expert VoIP & Infrastructures
% Classic CV template — Chef de Produit Téléphonie (fr)
% =============================================================================
\documentclass[11pt,a4paper,sans]{moderncv}
\moderncvstyle{classic}
\moderncvcolor{blue}
\usepackage[T1]{fontenc}
\usepackage{lmodern}
\usepackage{enumitem}
\usepackage{geometry}
\geometry{
  a4paper,
  top=0.7cm, bottom=0.6cm, left=1.1cm, right=1.1cm,
  scale=0.92
}
\setlength{\hintscolumnwidth}{3.2cm}

% === Informations Personnelles ===
\name{\cvfirstname}{{\cvlastname}}
\setlength{\makecvtitlenamewidth}{11cm}
\title{\bfseries Chef de Produit Téléphonie | Expert VoIP \& Infrastructures Télécoms}
\address{\cvlocation}{}
\phone[mobile]{\cvphone}
\email{\cvemail}
\social[linkedin]{\cvlinkedin}
\social[github]{\cvgithub}
\social[homepage]{\cvwebsite}

\nopagenumbers{}
% ── Personal data placeholders (overridden by ../shared/personal_data.tex) ────
\providecommand{\cvfirstname}{Firstname}
\providecommand{\cvlastname}{LASTNAME}
\providecommand{\cvemail}{user@example.com}
\providecommand{\cvphone}{+XX X XX XX XX XX}
\providecommand{\cvlocation}{City, Country}
\providecommand{\cvgithub}{github-username}
\providecommand{\cvlinkedin}{linkedin-username}
\providecommand{\cvwebsite}{website-url}
\input{../shared/personal_data}
\begin{document}
\makecvtitle
\vspace{-0.8cm}

% === PROFIL ===
\section{Profil}
\cvitem{}{Expert Télécoms avec \textbf{7 ans d'expérience} au sein du département ICT de la NIOC (opérateur interne d'envergure nationale). Spécialiste du déploiement d'infrastructures critiques : \textbf{VoIP/PABX, Réseaux TETRA, Faisceaux Hertziens, et Fibre Optique}. Profil hybride alliant maîtrise technique pointue (Mux/DeMux, BTS, SNMP) et gestion de projet opérationnelle. Master en Informatique \& DU Big Data obtenu (Montpellier).}

\vspace{-0.3cm}

% === COMPÉTENCES ===
\section{Compétences}
\cvitem{Gestion de Produit}{Pilotage d'offres télécoms, documentation technique/commerciale, rédaction de mémoires techniques, suivi de planning, harmonisation d'offres multi-sites.}
\cvitem{VoIP \& Téléphonie}{IP Centrex / PABX / PBX IP (Cisco, Avaya), SIP, QoS, Interconnexion de sites industriels.}
\cvitem{Infrastructures}{Réseaux TETRA (BTS, Antennes 4m), Faisceaux Hertziens (Microwave), Mux/DeMux, Fibre Optique, Câblage Cuivre, Antennes WiFi.}
\cvitem{Monitoring \& Automatisation}{SolarWinds (Expert), SNMP, Python (Scripting réseau), Dashboards temps réel.}
\cvitem{Gestion de Projet}{Supervision de chantiers, Contrôle Qualité (QA), Gestion des prestataires, Reporting management.}
\cvitem{Méthodologies}{Agile, Scrum, Cycle de vie produit (PLM), Rédaction de spécifications fonctionnelles.}

\vspace{-0.3cm}

% === EXPÉRIENCES ===
\section{Expériences Professionnelles}

\cventry{Jan 2025 -- Juin 2025}{\textbf{Ingénieur R\&D IA \& Automatisation (Stage)}}{toHero}{Montpellier}{}{
  Conception de pipelines Python pour l'extraction automatisée de données (PDF, Excel) via LLM / LangChain. Déploiement via Docker et intégration CI/CD.
}

\cventry{2021 -- 2024}{Ingénieur Support Télécoms \& Automatisation}{NIOC (ICT)}{Iran}{}{
  \begin{itemize}[leftmargin=*,label=$\circ$]
    \item \textbf{Opérateur Interne :} Gestion autonome de la connectivité nationale pour \textbf{1 500+ utilisateurs} sans recours à des prestataires externes.
    \item \textbf{VoIP \& Infrastructure :} Maintenance et optimisation du parc PABX Avaya/Cisco. Disponibilité garantie à \textbf{99,9\%}.
    \item \textbf{Automatisation :} Développement de scripts Python/SNMP pour le monitoring proactif, réduisant les interventions manuelles de \textbf{70\%}.
  \end{itemize}
}

\cventry{2017 -- 2019}{Ingénieur Support Réseaux (VoIP)}{NIOC (ICT)}{Iran}{}{
\begin{itemize}[leftmargin=*,label=$\circ$]
\item \textbf{Gestion de l'Offre VoIP :} Déploiement et configuration des solutions IP Centrex/PABX pour les sites industriels.
\item \textbf{Transition Numérique :} Modernisation des infrastructures legacy vers des solutions IP full-stack.
\end{itemize}
}

\cventry{2010 -- 2013}{Superviseur de Projet Télécom (TETRA) — Inspecteur QA}{NIOC}{Iran}{}{
  \begin{itemize}[leftmargin=*,label=$\circ$]
    \item \textbf{Déploiement National :} Supervision de l'installation du réseau TETRA sur \textbf{78 sites industriels}.
    \item \textbf{Infrastructure Radio :} Pilotage installation BTS, FH (Microwave) et systèmes Mux/DeMux.
    \item \textbf{Contrôle Qualité :} Validation de conformité, rapports de réception et rédaction des procédures.
  \end{itemize}
}

\vspace{-0.3cm}

% === FORMATIONS ===
\section{Formations}
\cventry{2025 -- 2026}{DU Big Data \& Data Science}{Université de Montpellier}{France}{}{Analyse de risques et modélisation prédictive via Python.}
\cventry{2024 -- 2025}{Master en Informatique (Bac+5)}{Université de Montpellier}{France}{}{Spécialisation Réseaux et IA -- Classé 1\textsuperscript{er} en Réseaux Avancés (16,33/20).}

\vspace{-0.2cm}

\section{Langues \& Certifications}
\cvitem{Langues}{Français (Avancé) | Anglais (Technique) | Persan (Maternel)}
\cvitem{Certifications}{\textbf{Formation technique certifiante — SIAE Microelettronica (Milan, Italie)} : Transmission Faisceaux Hertziens, Multiplexage (Mux/DeMux) et Systèmes NMS (2012).}

\end{document}
